% Generated draft: requires mathematical and source review.
\documentclass[11pt]{article}
\usepackage[margin=25mm]{geometry}
\usepackage{fontspec}
\setmainfont{Latin Modern Roman}
\usepackage{amsmath,amssymb,mathtools,unicode-math}
\setmathfont{Latin Modern Math}
\usepackage{enumitem,needspace,xurl,hyperref}
\setlength{\parindent}{0pt}
\setlength{\parskip}{5pt}
\setlength{\emergencystretch}{3em}
\newcommand{\N}{\mathbb{N}}
\newcommand{\Z}{\mathbb{Z}}
\newcommand{\Q}{\mathbb{Q}}
\newcommand{\R}{\mathbb{R}}
\newcommand{\C}{\mathbb{C}}
\newcommand{\F}{\mathbb{F}}
\newcommand{\A}{\mathbb{A}}
\newcommand{\PP}{\mathbb{P}}
\newcommand{\E}{\mathbb{E}}
\newcommand{\eps}{\varepsilon}
\newcommand{\dd}{\,\mathrm{d}}
\begin{document}
\section*{UnsolvedMath ID: 11000323}
\textbf{Source problem code:} catalogue entry
\subsection{Definitions and mathematical statement}
Retain the central extension (20), whose actual middle group and maps are not supplied, and write its abstract shape as
\[
 1\longrightarrow A\longrightarrow E\longrightarrow H_{g,1}\longrightarrow1,
 \qquad A=\Theta^3_{\Z}.
\]
Here \(H_{g,1}\) and \(\Theta^3_{\Z}\) are the source's groups, not defined in the excerpt. Centrality makes the coefficient action on \(A\) trivial.

\emph{Problem 11.2.} Study this extension through group cohomology and geometric topology, in particular determine its Euler, or central-extension, class
\[
 e(E)\in H^2(H_{g,1};\Theta^3_{\Z}).
\]
The extension itself is essential data: its class cannot be reconstructed merely from its kernel and quotient. The source writes this group as
\[
 \operatorname{Hom}(H_2(H_{g,1};\Z),\Theta^3_{\Z})
 \oplus\operatorname{Ext}_{\Z}^1(H_1(H_{g,1};\Z),\Theta^3_{\Z}).
\]
Under the usual universal-coefficient theorem there is an exact sequence with the Ext term injecting into \(H^2\) and the Hom term as quotient; a splitting is generally noncanonical. Thus a direct-sum coordinate description of the Euler class also requires a splitting choice and must not be treated as a canonical identification.

The source calls the problem extremely difficult and sketches a possible route to additive invariants of \(\Theta^3_{\Z}\), promising details in a then-forthcoming paper. Its trace construction yields classes
\[
 \widetilde t_{2k+1}\in
 H^2\!\left(\varprojlim_m\operatorname{Aut}_0(\Gamma/\Gamma_m)\right),
 \qquad k=1,2,\ldots,
\]
described as group versions of
\[
 t_{2k+1}\in H^2(\mathfrak h_{g,1}^{\Q})^{\operatorname{Sp}}
\]
from \(\S6\), Definition 6.1. It says a homomorphism \(\sigma\) is trivial on \(\Theta^3_{\Z}\), hence descends to
\[
 \bar\sigma:H_{g,1}\to
 \varprojlim_m\operatorname{Aut}_0(\Gamma/\Gamma_m).
\]
The groups, filtration \(\Gamma_m\), trace maps, cohomology coefficients, and the original \(\sigma\) are not defined here. No explicit formula identifying these trace classes with \(e(E)\), or an additive invariant, is supplied and none is invented.
\subsection{Short English statement}
Determine the central-extension class of the source's extension (20), and investigate its topology and possible additive invariants of \(\Theta^3_{\Z}\). The defining extension and trace setup are missing, and universal-coefficient direct-sum coordinates are generally noncanonical.
\subsection{Sources}
\begin{itemize}
\item Local UnsolvedMath registry, permanent ID 11000323, source code catalogue entry, Problem 11.2, PDF page 378. The chapter header identifies S. Morita; no individual proposer field is supplied.
\item Farb, \emph{Problems on Mapping Class Groups and Related Topics}. Supplied URL: \url{https://www.math.uchicago.edu/~farb/papers/mcgbook.pdf}; recorded access 2026-07-29; extraction pdf-font-structure. The edited volume contains 329 explicitly typeset headings across chapters, with repeated numbering disambiguated by page.
\item The essential extension (20), \(\S6\), Definition 6.1, and the then-forthcoming paper for details are referenced but not supplied or bibliographically identified.
\item Supplied status open and OPEN-TRIAGE. The summary says only ``Cannot verify.'' Current primary-literature status and redistribution terms for the publicly accessible author/publisher copy remain NEEDS\_REVIEW.
\end{itemize}
\end{document}
